\begingroup
\fontsize{12pt}{12pt}
\selectfont

\phantomsection
\pdfbookmark[1]{Задание на ВКР}{Задание на ВКР}

\vkrStandartHead 

\vspace{0.5cm}

\hfill%
\begin{minipage}{0.5\linewidth}
	\raggedleft
	\textbf{<<УТВЕРЖДАЮ>>}\\
	Заведующий кафедрой ИТиВС\\
	к.т.н., доц. Новоселова О.\,B.\\[3ex]
	\underline{\hspace{30mm}}\\
	<<\underline{\hspace{7mm}}>>.\underline{\hspace{30mm}} 2025 г.
\end{minipage}

\begin{center}
	\textbf{Задание}
	
	на выполнение выпускной квалификационной работы \\
	по направлению подготовки \\
	09.03.01 <<Информатика и вычислительная техника>>\\
	Профиль: <<Разработка программных комплексов в рамках цифровой трансформации предприятий>>
\end{center}

{
	\centering\fontsize{14pt}{12pt}\selectfont
	\begin{tabular} {l l l}
		Студент группы \vkrStudentGroup    & \hspace{0pt} &  Научный руководитель \\
		\vkrStudentShort			&  & 		\insertPersonInfo{Advisor} \vkrAdvisorShort\\
	\end{tabular}\\
}

\begin{center}
	\bfseries\fontsize{14pt}{12pt}\selectfont
	\parbox{0.9\linewidth}{\centering Тема: <<\vkrTheme>>}
	
\end{center}

\vfill

\noindent Тема утверждена приказом <<\underline{\hspace{7mm}}>>.\underline{\hspace{30mm}} 202\_ г. № \_\_\_\_

\noindent Срок сдачи ВКР на кафедру <<\underline{\hspace{7mm}}>>.\underline{\hspace{30mm}} 202\_ г.

\newpage
\fontsize{12pt}{18pt}
\selectfont
\begin{enumerate} [nosep,font=\bfseries, left=0pt] 
	\item \textbf{Описание задания на выполнение ВКР }
		\begin{enumerate} [nosep, label*=\arabic*.,font=\bfseries]
			\item \textbf{Тип ВКР ---} исследовательская работа
			\item \textbf{Цель исследования ---} ...
			\item \textbf{Объект исследования --- } ...
			\item \textbf{Предмет исследования ---} ...
			\item \textbf{Методы исследования --- } ...
			\item \textbf{Задачи исследования:}
				\begin{enumerate} [nosep, label*=\arabic*.]
					\item Задача 1
					\item Задача 2
					\item Задача 3
					\item Задача 4
				\end{enumerate}
		\end{enumerate}
	\item \textbf{Требования к выполнению ВКР}
		\begin{enumerate} [nosep, label*=\arabic*.,font=\bfseries]
			\item \textbf{Соблюдение требований законодательной базы и стандартов}
				\begin{enumerate} [nosep, label*=\arabic*.]
					\item  Образовательная программа высшего образования в бакалавриате ФГБОУ ВО «МГТУ «СТАНКИН» по направлению подготовки 09.03(04).01(04)  «Направление» для направленности (профиля) подготовки «Профиль » (утв. 07.04.2016 (изменена и дополнена) .
					\item Внутренний нормативный документ. П 01-04/264/2017. Положение о государственной итоговой аттестации по образовательным программам высшего образования – программам бакалавриата, программам специалитета и программам магистратуры: [утверждено приказом ректора от 31.08.2017 г. №431/1, одобрено решением ученого совета Университета от 31.08.2017 г., протокол № 10/17] – ФГБОУ ВО «МГТУ «СТАНКИН». 
					\item Внутренний нормативный документ. П 01-04/438/2021. Положение о выпускной квалификационной работе обучающихся по образовательным программам высшего образования – программам бакалавриата, программам специалитета, программам магистратуры: [утверждено приказом врио ректора от 30.03.2021 г. №177/1, одобрено решением ученого совета Университета от 25.12.2020 г., протокол № 11/20] – ФГБОУ ВО «МГТУ «СТАНКИН».
				\end{enumerate}
			\item \textbf{Дополнительные требования}
				\begin{enumerate} [nosep, label*=\arabic*.]
					\item 	Оформление раздела «Список литературы» по национальному стандарту ГОСТ Р 7.0.100-2018 “Библиографическая запись. Библиографическое описание. Общие требования и правила составления”
					\item Результаты исследования должны быть опубликованы в виде научных статей и тезисов докладов (не менее 1 ), а также доложены на научно-технических конференциях и семинарах. 
				\end{enumerate}	
		\end{enumerate}

	\item \textbf{Срок сдачи оформленной квалификационной работы на кафедру ---} вторая декада мая 2023~г. 
		
\end{enumerate}

\vspace{2ex}%
\noindent%
\begin{tabular}{   @{}p{0.75\linewidth}@{}p{0.25\linewidth}@{}    }
	Исполнитель 												&  \vkrStudentShort   \\[2ex]
	
	Научный руководитель\\\insertPersonInfo{Advisor} каф. ИТиВС			&   \vkrAdvisorShort 
\end{tabular}


\endgroup